\paragraph{A uniform length bound for odd cocycles}
In this paragraph write \(N=N_{\mathrm{cyc}}\) and \(y=N+t+3\). The module is zero if \(\nu=0\), so take \(\nu\ge1\). At exponent \(a=0\), global class field theory
identifies the everywhere unramified continuous homomorphisms \(G_E\to\mathcal O/\ell^\nu\)
with the homomorphisms from \(\operatorname{Cl}(E)\). A cyclic factor of order \(n_j\)
contributes length at most \(v_\ell(n_j)\), so the length is at most \(v_\ell(h_E)\).

Suppose \(0<a<t\). The coefficients of the Bernoulli polynomials give the integers
\(B_j=(a+2)!\,N^{a+1}B_{a+1}(j/N)\).
Define \(g(j)=2B_j+(-1)^a\mathbf 1_{j=1}+\mathbf 1_{j=-1}\) using the least nonnegative
representatives modulo \(N\). The matrix \(T_{s,j}=g(j/s)\) on the units modulo \(N\) has
eigenvalue \(2\sum_j\psi(j)B_j+(-1)^a+\psi(-1)\) at \(\psi\).
When \(\psi(-1)=-(-1)^a\), the Bernoulli sum is a nonzero multiple of \(L(-a,\psi)\)
by the functional equation and the nonvanishing Euler product in \(\operatorname{Re}s>1\).
The Euler factors from passing to the inducing primitive character are nonzero since \(a>0\).
For the other parity the Bernoulli sum vanishes and the eigenvalue is \(2(-1)^a\). Take
\[
 X_a=|\det T|\prod_{c\in\{1+N,1+2N\}}(c^a-1)(c^{a+1}-1).
\]
The entries have absolute values at most \(64(a+2)!N^{a+1}(a+1)!\).
Expanding the determinant gives \(0<X_a\le2^{(N+a+3)^6}\).
At least one of the two \(c\)'s is prime to any given \(\ell\).

Let \(c\equiv1\pmod N\) be prime to \(\ell\) and put \(R=N\ell^\nu\).
If \(\psi(-1)=-(-1)^a\), in \(\mathcal O/\ell^\nu\) the following sum divides
\((c^{a+1}-1)X_a\).
\[
 S_c=\sum_{j\in(\mathbb Z/R)^\times}
 \left(\left\lfloor\frac{cj}{R}\right\rfloor
       -\left\lfloor\frac{c[-j]_R}{R}\right\rfloor\right)j^a\psi(j).
\]
Here \(j\) and \([-j]_R\) are least nonnegative representatives.
Applying the adjugate to the character vector and taking its coordinate at \(1\) shows
that the eigenvalue \(2\sum_j\psi(j)B_j\) divides \(X_a\).
Voronoi's congruence writes its product with \(c^{a+1}-1\) as a multiple of a sum of floors.
Comparing the full sums at \(N\ell^\nu\) and \(N\ell^{\nu-1}\) removes the terms indexed
by multiples of \(\ell\) and introduces the unit \(1-\ell^a\psi(\ell)\) on the left.
Replacing \(j\) by \(-j\) in the sum over the units gives the difference of floors.
\leanlinks{\href{\upstream/grh/src/Lib/Cyclotomic/Voronoi.lean\#L393}{Construction of the integer from Bernoulli polynomials}
  \leanlinksep \href{\upstream/grh/src/Lib/Cyclotomic/Kolyvagin/ReciprocityComponentBounds.lean\#L576}{The congruence restricted to unit residues}}

Choose lifts \(\sigma\) from \(C_N=(\mathbb Z/N)^\times\) to \(G_{\mathbb Q}\) with \(\sigma(1)=1\).
For complex conjugation \(c_\infty\), put
\[
 B(u,h)=u(c_\infty h c_\infty^{-1}h^{-1})u^{-1},\qquad
 z_\psi(h)=\sum_{v\in C_N}(\epsilon_{\mathrm{cyc}}^a\psi)(\sigma(v))\,z(B(\sigma(v),h)).
\]
The cyclotomic quotient is abelian, so \(B(u,h)\in G_E\).
We have \(B(u,hh')=B(u,h)B(uh,h')\), and the cyclotomic character is trivial on these commutators.
Left multiplication of \(u\) by an element of \(G_{\mathbb Q(\mu_N)}\) changes \(z(B(u,h))\)
by the conjugation factor that cancels the change in \(\epsilon_{\mathrm{cyc}}^a(u)\).
Reindexing by \(v\mapsto v\overline h\) gives the cocycle law for \((\epsilon_{\mathrm{cyc}}^a\psi)^{-1}\).
The commutator at \(c_\infty\) is trivial, and at the prescribed local elements \(z\) vanishes
on both conjugates that form the commutator. Character orthogonality gives, for \(h\in G_E\),
\[
 \sum_\psi z_\psi(h)=\varphi(N)\bigl(z(c_\infty h c_\infty^{-1})-z(h)\bigr).
\]
An even \(\epsilon_{\mathrm{cyc}}^a\psi\) gives a component invariant under conjugation by \(c_\infty\).
If the odd components vanish, comparing the formula at \(h\) and \(c_\infty h c_\infty^{-1}\)
and using the oddness of \(z\) yields \(4\varphi(N)z=0\).
\leanlinks{\href{\upstream/grh/src/Lib/Cyclotomic/OddSha.lean\#L414}{The character components of a cocycle}}

If an odd component has \(4X_a^2z_\psi\ne0\), then \(2X_a^2z_\psi\) is nonzero on
\(G_{\mathbb Q(\mu_R)}\). Indeed, vanishing on this subgroup would give descent to an abelian
quotient, where the value zero and the action \(-1\) at \(c_\infty\) would force twice the cocycle
to vanish. By Chebotarev choose a Frobenius \(\sigma_p\) at \(p\equiv1\pmod R\) where it is nonzero.
Take \(c\in\{1+N,1+2N\}\) prime to \(\ell\).
The Jacobi sum \(J(1/R)\) is a unit away from \(p\), and the valuations of its conjugates
above \(p\) are the differences of floors in \(S_c\).
Artin reciprocity for this element gives \(S_c\,2(c^a-1)z_\psi(\sigma_p)=0\).
Using \(S_c\mid(c^{a+1}-1)X_a\) and \((c^a-1)(c^{a+1}-1)\mid X_a\) in turn yields
\(2X_a^2z_\psi(\sigma_p)=0\), a contradiction.
Thus \(Y_a=4\varphi(N)\cdot4X_a^2\) annihilates the original odd cocycles.
\leanlinks{\href{\upstream/grh/src/Lib/Cyclotomic/Kolyvagin/ReciprocityComponentBounds.lean\#L670}{Annihilating each component using the Jacobi sums}}

We next find at most \(y^{20}\) evaluation points to detect a cocycle.
For \(\ell\le y^2\), put \(D=N\ell^2\) and \(K=\mathbb Q(\mu_D)\).
Use representatives for \(G_E/G_K\) and lifts of generators of the abelian quotient of \(G_K\)
unramified outside \(\ell\) and killed by \(D\).
The order of the latter quotient divides \(h_K\) times the index of local \(D\)-th powers above
\(\ell\). The cyclotomic discriminant, Minkowski's class number bound and the local index bound
give \(2^{D^3+5D^2+D}\) as an upper bound for this product.
Cyclic decomposition gives at most \(\varphi(D)+D^3+5D^2+D\le D^4\le y^{20}\) points.
For \(z\) vanishing at these points, induct on the exponent of a power of \(\ell\) killing \(z\).
The cocycle \(Dz\) vanishes at the same points and is killed by a smaller power, so \(Dz=0\).
Then \(G_K\) acts trivially on the values and \(z|_{G_K}\) is a homomorphism through the indicated
abelian quotient. Its vanishing on generators and coset representatives implies \(z=0\).
\leanlinks{\href{\upstream/grh/src/Lib/Cyclotomic/OddSha.lean\#L91}{The class number, the local index and the evaluation points at small primes}}

For \(\ell>y^2\) the Kolyvagin derivatives of Gauss sums give
\(m_a=\operatorname{ord}_\ell(X_a)+1\) evaluation points per odd component.
Suppose no \(m_a\) points detect the component. For a set \(T\) of at most \(m_a\) primes congruent
to \(1\) modulo \(N\ell\), choose a component vanishing at one Frobenius above each prime of \(T\)
and take its last nonzero \(\ell\)-power multiple. The resulting \(w\) has \(\ell w=0\) and
is trivial above \(T\): the action on the \(\ell\)-torsion coefficients is trivial at these primes,
and every decomposition group is topologically generated by inertia and the conjugate Frobenius.

Choose \(c\equiv1\pmod N\) to be a primitive root modulo \(\ell\), so \(c^a-1,c^{a+1}-1\) are units.
Choose an inertia element \(\gamma\) above \(\ell\) at which the dual character
\(\epsilon_{\mathrm{cyc}}^{a+1}\psi\) takes the value \(c^{a+1}\).
To construct a finite set \(Q\) of auxiliary primes congruent to \(1\) modulo \(cN\ell^{m_a}\),
consider a dual cocycle unramified outside \(\ell\) over \(\mathbb Q(\mu_N)\) and nonzero on
\(G_{\mathbb Q(\mu_{cN\ell^{m_a}})}\), paired with a nonzero odd component. The subgroup is no
union of the kernels of the two restrictions, since both are proper there.
Chebotarev gives a prime at whose Frobenius both are nonzero. Use one such prime for each pair.
There are finitely many cocycles since restriction to a subgroup where the coefficient action is
trivial factors through a finite class field theory quotient and only finitely many values
are possible on coset representatives. Next choose \(p\equiv1\pmod {N\ell^{m_a}\prod_{r\in Q}r}\)
above which a nonzero \(\ell\)-torsion component is nontrivial.

Consider the following property of \(T\). Every dual cocycle modulo \(\ell^{m_a}\) killed by
\(\ell\) and unramified over \(\mathbb Q(\mu_N)\) outside \(\ell,T\) is unramified outside \(\ell\)
over that field. The empty set has the property. If \(w\) is trivial above \(T\) but nontrivial
above \(r\equiv1\pmod{N\ell}\), then \(T\cup\{r\}\) also has it. To see this, apply equally to
each place above \(r\) the following calculation of the inertia value \(z=\kappa(\tau)\)
of any such dual cocycle \(\kappa\). Extend the coordinates in an \(\mathbb F_\ell\)-basis
\((e_i)\) of \(\mathcal Oz\) to homomorphisms to \(\mu_\ell\).
Over \(\mathbb Q(\mu_{N\ell})\), Hilbert 90 applied to these coordinates gives
Kummer elements \(\alpha_i\) with \(\alpha_i^\ell\) in that field.
Write \(m_{i,u}\) for the orders of its conjugates above \(r\), and
\(\tau(r^{1/\ell})/r^{1/\ell}=\zeta_\ell^j\).
The element \((r^{m_{i,u}/\ell}/\sigma(u)\alpha_i)^\ell\) is a local unit.
Inertia therefore fixes \(r^{m_{i,u}/\ell}/\sigma(u)\alpha_i\) up to a principal unit.
Comparing the exponents of \(\zeta_\ell\) gives
\(u^a\psi(u)\sum_i m_{i,u}j e_i=z\), where the factor \(u\) from conjugating
\(\mu_\ell\) cancels one power of the dual character.
Artin reciprocity for each \(\alpha_i^\ell\) shows that
\(\sum_u m_{i,u}u^a\psi(u)\) annihilates \(2(c^a-1)w(\operatorname{Frob}_r)\ne0\).
Each of these sums lies in the maximal ideal of \(\mathcal O\).
Summing the coordinate identities over \(u\) now puts
\(\varphi(N\ell)z\) in the maximal ideal times \(\mathcal Oz\).
Since \(\varphi(N\ell)\) is a unit, \(z=0\). In particular \(\{p\}\) has the property.
\leanlinks{\href{\upstream/grh/src/Lib/Cyclotomic/Kolyvagin/ReciprocityComponentBounds.lean\#L430}{The Kummer calculation in reciprocity modulo \(\ell\)}}

For \(S\subset Q\), write \(\kappa_S\) for the Gauss-sum derivative cocycles for \(p,Q\).
To construct them, take \(\ell^{m_a}\)-th roots \(Y(x)\) of the Jacobi family \(J(x)\).
For \(g\in G_{\mathbb Q}\) take the exponent in \(\mu_{\ell^{m_a}}\) of
\(\gamma(Y(x))/(g\gamma g^{-1})(Y(x))\) on arguments killed by \(N\ell^{m_a}\prod_{r\in S}r\).
For each \(r\in S\) apply the operator taking \(f\) to
\(\sum_{u\in(\mathbb Z/r)^\times}\log_r(u)f(x+u/r)\), where \(\log_r\) is onto
\(\mathbb Z/\ell^{m_a}\). Then sum at \(x=j/(N\ell^{m_a})\) with weights \(j^a\psi(j)\).
The resulting cocycles vanish at \(\gamma\) and are unramified over \(\mathbb Q(\mu_N)\)
outside \(\ell,p,S\). For \(r\in Q\setminus S\), choose an inertia element acting on \(\mu_r\)
as \(v^{-1}\) with \(\log_r(v)=1\). Substitution by \(v\) shifts the weight \(\log_r\) by \(1\).
The distribution relation of the Jacobi family and reduction at \(r\) identify the resulting
value of \(\kappa_{S\cup\{r\}}\) on inertia with \(\kappa_S(\operatorname{Frob}_r)\).
Here roots of unity of order prime to \(r\) are distinguished by reduction.
\leanlinks{\href{\upstream/grh/src/Lib/Cyclotomic/Kolyvagin/DerivativeCocycles.lean\#L922}{The Frobenius and inertia calculation for the derivatives}}

When \(\#S\le i\) and \(\{p\}\cup S\) has the property above, show in descending order from \(i=m_a\)
that \(\ell^i\kappa_S\) is unramified outside \(\ell\) over \(\mathbb Q(\mu_N)\).
Initially \(\ell^{m_a}\kappa_S=0\). At stage \(i\), if \(u=\ell^{i+1}\kappa_S=0\), the property
of the set applies. If \(u\ne0\), the value zero at \(\gamma\) and the unit difference of the dual
character from \(1\) imply that \(u\) is nonzero on the subgroup over the specified cyclotomic
field. Pair it with a nonzero \(w\) trivial above \(\{p\}\cup S\), and choose \(r\) from \(Q\)
where both Frobenius values are nonzero. Then \(r\notin S\), and adjoining it preserves the property.
By induction \(\ell^{i+1}\kappa_{S\cup\{r\}}\) is unramified above \(r\), so the equality on inertia
contradicts \(u(\operatorname{Frob}_r)\ne0\). Hence \(\kappa_\varnothing\) is unramified above \(p\).
Its inertia value formula yields \((\epsilon_{\mathrm{cyc}}^{a+1}\psi(\gamma)-1)S_c=0\).
The first factor and \(c^{a+1}-1\) are units, so the Stickelberger congruence gives
\(\ell^{m_a}\mid X_a\), contradicting the definition of \(m_a\).
\leanlinks{\href{\upstream/grh/src/Lib/Cyclotomic/Kolyvagin/ReciprocityComponentBounds.lean\#L1029}{The descending induction}}

If \(\gamma_{\psi,i}\) are the points for the components, use \(B(\sigma(v),\gamma_{\psi,i})\)
for \(z\). Vanishing at these points kills every odd component, and \(4\varphi(N)<\ell\) gives
\(z=0\). The number of points is at most \(\varphi(N)^2m_a\le y^{20}\).
For small and large primes alike, evaluation therefore embeds the module into at most
\(y^{20}\) copies of \((\mathcal O/\ell^\nu)[Y_a]\), each of length at most \(v_\ell(Y_a)\).
For all exponents and primes we can take \(C_{\mathrm{odd}}=h_E(\prod_{0<a<t}Y_a)^{y^{20}}\).
There are at most \(y\) factors and each \(Y_a\le 2^{4y^6}\), so
\(C_{\mathrm{odd}}\le h_E2^{4y^{27}}\).
\leanlinks{\href{\upstream/grh/src/Lib/Cyclotomic/OddSha.lean\#L507}{The uniform bound for all exponents and primes}}
